\chapter{Camada TTY, syscalls e getty}
\label{cap:tty}

\section{Do renderizador VGA aos consoles virtuais}

At\'e aqui o driver VGA era dono tanto do conte\'udo da tela quanto da posi\c{c}\~ao
do cursor. Para suportar v\'arios terminais como no Linux, o c\'odigo foi
dividido em duas camadas:

\begin{itemize}
  \item \code{vga.rs} \'e um \emph{renderizador}. Ele conhece a matriz de
        c\'elulas 80$\times$25 em \code{0xB8000}, os registradores CRTC do
        cursor e nada sobre linhas de comando ou entrada;
  \item \code{tty.rs} mant\'em oito consoles virtuais, \code{tty0} a
        \code{tty7}. Cada console tem seu pr\'oprio buffer de c\'elulas,
        cursor e fila de entrada.
\end{itemize}

Somente o console ativo \'e copiado para o framebuffer de hardware. Trocar de
console \'e, portanto, uma c\'opia de buffer mais a atualiza\c{c}\~ao do cursor,
a mesma id\'eia usada pela camada de terminal virtual (\code{vt}) do Linux.

\begin{lstlisting}[caption={Estado de cada console em \code{tty.rs}.},label={lst:ttytty}]
#[derive(Clone, Copy)]
struct Tty {
    cells: [u16; vga::CELLS],
    row: usize,
    column: usize,
    input: [u8; INPUT_SIZE],
    head: usize,
    tail: usize,
}

static TTYS: SpinLock<[Tty; TTY_COUNT]> =
    SpinLock::new([Tty::new(); TTY_COUNT]);
static ACTIVE: AtomicUsize = AtomicUsize::new(0);
\end{lstlisting}

A fila de entrada \'e um pequeno buffer circular de 256 bytes. O
\emph{produtor} \'e \code{pump\_input}, que drena o teclado PS/2 e a COM1; o
\emph{consumidor} \'e o shell ou uma chamada de sistema lendo do console ativo.

\section{Troca de console}

O driver de teclado agora exp\~oe dois tipos de resultado:

\begin{lstlisting}[caption={Tipo de resultado do teclado.},label={lst:keyenum}]
pub enum Key {
    Char(u8),
    Console(usize),
}
\end{lstlisting}

O driver acompanha Shift e Alt. Com Alt pressionado, os scancodes
\code{0x3B}..\code{0x42} (F1..F8) viram \code{Key::Console(0..7)}. A camada
TTY traduz isso em \code{switch(index)} e redesenha o buffer escolhido. A
mesma opera\c{c}\~ao est\'a dispon\'ivel no shell:

\begin{lstlisting}[style=shell,caption={Trocando de console pelo shell.},label={lst:ttyshell}]
rinux> tty
console ativo: tty0 (disponiveis: tty0..tty7)
rinux> tty 3
rinux> tty
console ativo: tty3 (disponiveis: tty0..tty7)
\end{lstlisting}

Como o shell \'e um \'unico la\c{c}o s\'incrono, trocar de console durante a
edi\c{c}\~ao de uma linha descarta a linha parcial e imprime um novo prompt no
console novo. Essa \'e uma simplifica\c{c}\~ao deliberada: ainda n\~ao existe
um \code{task\_struct} por terminal.

\section{Sequ\^encias de escape ANSI/VT100}

A sa\'ida n\~ao se limita a caracteres imprim\'iveis. Cada \code{Tty} mant\'em
sua pr\'opria m\'aquina de estados, para que um programa use as mesmas
sequ\^encias de escape de um terminal Linux sem saber qual console est\'a ativo:

\begin{lstlisting}[caption={Estado de escape guardado em cada console.},label={lst:ttyescape}]
enum EscapeState {
    Ground, // texto normal
    Start,  // recebeu ESC
    Csi,    // recebeu ESC [ e le parametros
}

struct Tty {
    // ...
    escape: EscapeState,
    csi_params: [u16; 4],
    csi_count: usize,
    csi_value: u16,
    csi_digits: bool,
}
\end{lstlisting}

\code{write\_byte} encaminha \code{0x1B} para a m\'aquina de estados. No estado
CSI o interpretador aceita par\^ametros decimais separados por \code{;} e um
byte final; marcadores privados como \code{?} e bytes intermedi\'arios s\~ao
consumidos sem serem impressos:

\begin{lstlisting}[caption={Despacho de ESC e CSI em \code{tty.rs}.},label={lst:ttywritebyte}]
fn write_byte(&mut self, byte: u8) {
    match self.escape {
        EscapeState::Ground => {
            if byte == 0x1B {
                self.escape = EscapeState::Start;
            } else {
                self.write_plain(byte);
            }
        }
        EscapeState::Start => {
            if byte == b'[' {
                self.escape = EscapeState::Csi;
                self.csi_params = [0; 4];
                self.csi_count = 0;
                // ...
            } else {
                self.escape = EscapeState::Ground;
            }
        }
        EscapeState::Csi => self.write_csi(byte),
    }
}
\end{lstlisting}

\begin{table}[H]
\centering
\caption{Sequ\^encias CSI suportadas.}
\label{tab:ansi}
\begin{tabularx}{\linewidth}{@{}lX@{}}
\toprule
\textbf{Sequ\^encia} & \textbf{Efeito} \\
\midrule
\code{ESC[2J}, \code{ESC[3J} & Apaga a tela inteira. \\
\code{ESC[J} & Apaga do cursor at\'e o fim da tela. \\
\code{ESC[K} & Apaga do cursor at\'e o fim da linha. \\
\code{ESC[H}, \code{ESC[<linha>;<coluna>H} & Leva o cursor ao in\'icio ou a
uma posi\c{c}\~ao base 1, limitada a 80$\times$25. \\
\code{ESC[A}, \code{ESC[B}, \code{ESC[C}, \code{ESC[D} & Move o cursor para
cima, baixo, direita ou esquerda (1 por padr\~ao, ou a quantidade dada). \\
\code{ESC[...m} & SGR \'e interpretado e ignorado; cores ainda n\~ao existem. \\
\bottomrule
\end{tabularx}
\end{table}

O comando \code{clear} agora \'e sa\'ida comum, n\~ao um caso especial do
renderizador:

\begin{lstlisting}[caption={Implementa\c{c}\~ao do comando \code{clear}.},label={lst:ttyclear}]
pub fn clear() {
    write_bytes(b"\x1b[2J\x1b[H");
}
\end{lstlisting}

\code{ESC[2J} apaga as c\'elulas mas, seguindo o comportamento VT100, n\~ao
move o cursor; a segunda sequ\^encia \code{ESC[H} o leva ao canto superior
esquerdo. Como \code{clear} passa por \code{write\_bytes}, a sequ\^encia
tamb\'em \'e espelhada na COM1, de modo que um terminal serial tamb\'em limpa.
Cores SGR, tela alternativa, scrollback e visibilidade do cursor ainda est\~ao
fora do escopo.

\section{ABI das chamadas de sistema}

A porta \code{0x80} j\'a existia para os demos em ring 3. Agora ela despacha
para \code{syscall.rs}, com a conven\c{c}\~ao usual de registradores:

\begin{table}[H]
\centering
\caption{Chamadas de sistema implementadas.}
\label{tab:syscalls}
\begin{tabularx}{\linewidth}{@{}clX@{}}
\toprule
\textbf{N\'umero} & \textbf{Nome} & \textbf{Argumentos e resultado} \\
\midrule
0 & \code{exit} & \code{rdi} = status; n\~ao retorna. \\
1 & \code{write} & \code{rdi} = buffer, \code{rsi} = tamanho;
devolve o n\'umero de bytes. \\
2 & \code{read} & \code{rdi} = buffer, \code{rsi} = tamanho,
\code{rdx} = flags (bit 0 desliga o eco); devolve os bytes lidos. \\
3 & \code{get\_tty} & sem argumentos; devolve o \'indice do console ativo. \\
\bottomrule
\end{tabularx}
\end{table}

\subsection{Valida\c{c}\~ao de ponteiros do usu\'ario}

Uma chamada de sistema roda em ring 0 com as mesmas tabelas de p\'agina do
usu\'ario. Tocar em um endere\c{c}o de usu\'ario n\~ao mapeado n\~ao geraria
um page fault de usu\'ario; seria um acesso do kernel. Por isso \code{write}
e \code{read} chamam \code{paging::ensure\_user\_range} antes de
desreferenciar o buffer:

\begin{lstlisting}[caption={Valida\c{c}\~ao e mapeamento sob demanda de um buffer do usu\'ario.},label={lst:ensureuser}]
pub fn ensure_user_range(addr: u64, len: u64) -> Result<(), &'static str> {
    let end = addr.checked_add(len).ok_or("intervalo invalido")?;
    if addr < USER_BASE || end > USER_BASE + USER_REGION_SIZE {
        return Err("ponteiro fora da regiao do usuario");
    }
    // Mapeia com U/S = 1 cada pagina ausente de [addr, end).
    // ...
    Ok(())
}
\end{lstlisting}

A fun\c{c}\~ao rejeita intervalos fora da regi\~ao do usu\'ario e materializa
cada p\'agina ausente com \code{PAGE\_USER | PAGE\_WRITE}. \'E o equivalente
did\'atico dos cuidadosos auxiliares \code{copy\_from\_user}/\code{copy\_to\_user}
do Linux.

\subsection{Leitura bloqueante}

A chamada \code{read} \'e s\'incrona. Enquanto n\~ao h\'a byte dispon\'ivel,
ela chama \code{tty::pump\_input} e \code{spin\_loop}. A porta de interrup\c{c}\~ao
limpa \code{IF}, portanto o timer n\~ao avança durante a espera na
implementa\c{c}\~ao atual. Bytes imprim\'iveis s\~ao gravados no buffer do
usu\'ario e ecoados quando a flag de sem-eco est\'a limpa; Backspace remove o
\'ultimo byte; Enter termina a linha e devolve o caractere de nova linha.

\section{getty e login em ring 3}

O primeiro programa de usu\'ario de verdade \'e um \emph{getty}/login. Ele \'e
escrito em assembly dentro de \code{user.rs}, em uma se\c{c}\~ao pr\'opria, e
copiado como um blob independente de posi\c{c}\~ao para
\code{paging::USER\_CODE}:

\begin{lstlisting}[caption={Fluxo de login do programa em ring 3.},label={lst:gettyasm}]
getty_start:
    GETTY_WRITE getty_login_msg, getty_login_end   # write(2)
    subq $128, %rsp                                # buffer na pilha
    movq %rsp, %rdi                                # buf
    movq $32, %rsi                                 # len
    xorq %rdx, %rdx                                # echo
    movq $2, %rax                                  # read
    int $0x80
    # compara "root" e pede "Password: "
    # segundo read com flags = 1 (sem eco)
    # compara "rinux" e escreve a mensagem final
    # exit(0) ou exit(1)
\end{lstlisting}

Usar uma se\c{c}\~ao dedicada com c\'odigo e strings mant\'em todas as
refer\^encias RIP-relativas v\'alidas depois da c\'opia. O buffer de leitura
fica na pilha do usu\'ario, cuja p\'agina \'e mapeada sob demanda pela primeira
chamada de sistema que a toca.

As credenciais da demonstra\c{c}\~ao s\~ao:

\begin{lstlisting}[style=shell,caption={Login bem-sucedido.},label={lst:gettyout}]
rinux> getty
[GETTY] iniciando login em tty0...
rinux login: root
Password:
Bem-vindo ao rinux, root!
[GETTY] login concluido, voltando ao shell
\end{lstlisting}

\code{getty n} primeiro troca para o console \code{n} e depois executa o
programa nele. Em um teste automatizado pela serial, as credenciais devem ser
enviadas somente depois que o prompt de login daquele console aparecer, pois a
porta serial \'e um \'unico fluxo global.

\section{Rela\c{c}\~ao com o Linux}

\begin{table}[H]
\centering
\caption{Correspond\^encia com o subsistema TTY do Linux.}
\label{tab:ttflinux}
\begin{tabularx}{\linewidth}{@{}lX@{}}
\toprule
\textbf{Linux} & \textbf{Aqui} \\
\midrule
\code{drivers/tty/tty\_io.c} & \code{src/tty.rs}: filas de leitura/escrita. \\
\code{drivers/tty/vt/vt.c} & \code{src/tty.rs}: consoles virtuais e troca. \\
\code{drivers/video/console/vgacon.c} & \code{src/vga.rs}: renderizador de
c\'elulas. \\
\code{drivers/tty/n\_tty.c} & N\~ao implementado; as linhas s\~ao tratadas
diretamente. \\
\code{arch/x86/entry/syscalls/syscall\_64.tbl} & \code{src/syscall.rs}. \\
\code{init/getty.c}, \code{login} & Blob ring 3 em \code{src/user.rs}. \\
\bottomrule
\end{tabularx}
\end{table}

\section{Limita\c{c}\~oes atuais}

\begin{itemize}
  \item Existe um shell, n\~ao um shell por terminal. Trocar de console move o
        mesmo la\c{c}o para outro buffer.
  \item O \code{getty} n\~ao \'e residente: ele roda uma vez, de forma
        s\'incrona, e volta ao shell. Um sistema real mant\'em um
        \code{getty} esperando em cada TTY.
  \item A serial \'e um \'unico fluxo; o teclado \'e o que realmente distingue
        os consoles no uso interativo.
  \item N\~ao h\'a disciplina de linha: nem modo can\^onico, sinais, controle
        de jobs, \code{Ctrl+C}, \code{Ctrl+D} ou ioctls de terminal.
  \item \code{read} espera com interrup\c{c}\~oes desligadas; o timer congela
        durante um prompt de login.
\end{itemize}
